%
% File: chap04.tex
% Author: Solal Rapaport
% Description: History Alterations in Public Git Repositories (based on
% contributions/altered_histories.pdf). The methodology reuses the shared
% snapshot-pair stance established in Chapter~\ref{chap:tags} and details only
% the differences.
%
\chapter{History Alterations in Public Git Repositories}
\label{chap:histories}

\section{Introduction}
\label{sec:hist-intro}

Chapter~\ref{chap:tags} established that tags---the named references developers use to designate software releases---are mutable in practice.
Yet a tag can remain perfectly stable while the commits it ultimately reaches through its ancestry have been silently replaced: mutable release labels are only one surface of repository mutability, and this chapter moves the investigation inward, to the commit history that public branches advertise at each archival observation.

For the most part, Git repositories are expected to be append-only: developers add \emph{new} commits and exchange them by pushing to and pulling from public repositories~\cite{ProGit2014}.
It is nonetheless technically possible to \emph{alter the previously recorded history} of a repository~\cite[Chapter 7.6, Rewriting History]{ProGit2014}, typically by ``amending'' a recent commit to change its content or metadata, or by ``rebasing'' a portion of history to merge, split, or linearize commits.
Downstream recipients have no way to tell, from the repository alone, that this happened.
That silence is what turns history alteration into a supply-chain concern rather than a curiosity.
History rewrites are useful to polish local or work-in-progress branches before sharing them, and are even the default behaviour in some workflows~\cite{DBLP:journals/infsof/RiosEE22}: pull-request branches, for example, are often rewritten to keep only a single commit that can be merged when the pull request is ready.
The same capability becomes a liability when exercised on public branches that downstream recipients already track.
Here the stakes differ from a moved tag: a retargeted tag misdirects \emph{which} commit is fetched, whereas a rewritten history changes \emph{what} that commit---and every commit after it---actually contains.
Two converging reasons make this matter.
First, it breaks the pull/push workflow~\cite{DBLP:conf/issre/JiCYM20}, which relies on the stability of commit identifiers as checkpoints for history exchanges.
Second, and more consequentially, by breaking commit identities it opens the door to supply-chain attacks: together with a legitimate-looking history rewrite, an attacker can sneak in malicious changes that are difficult to audit because of the confusion caused by commit-identity modifications~\cite{ladisa2023supplychain,torresarias2016gitmetadata}.

Studying history alterations empirically is difficult for a structural reason: a rewrite, by construction, destroys the evidence needed to detect it---once the prior history is gone from the repository, there is nothing left there to compare the new state against.
Consequently, little was known, before this study, about the amount, characteristics, and impact of history alterations in public code.
As in the previous chapter, Software Heritage solves this methodological problem: successive crawls retain prior object identities even when forges overwrite refs~\cite{swhcacm2018,cise-2020-doi}.
That archival property makes it possible to ask, at ecosystem scale, how often public histories change in ways that remove or replace commits, how those events distribute across branches and projects, and whether the file-level footprints resemble benign hygiene (for example, cleaning credentials) or governance-sensitive patterns (for example, retroactive license edits).

\paragraph{What this chapter establishes.}
Tag-level evidence cannot settle this question.
A project can show perfectly stable release tags while the commit content and metadata those tags reach through are extensively rewritten.
The risk this chapter targets is therefore not reducible to tag mutability alone.
This chapter extends the study we published at ASE~2025~\cite{DBLP:conf/kbse/RapaportPTZ25} along five complementary lines that quantify that risk directly.
It quantifies how often publicly archived Git histories change between successive observations and how those events distribute across repositories (\Cref{subsec:hist-freq}); identifies which branches carry the largest share of alteration events and how the likelihood of rewriting relates to repository visibility on GitHub (\Cref{subsec:hist-dist-branches,subsec:hist-affected-projects}); applies a taxonomy that separates metadata-only rewrites from directory and file changes (\Cref{sec:hist-taxonomy}); examines two recurrent file-level patterns in depth---retroactive license changes and removals of files whose names suggested secrets (\Cref{sec:hist-security}); and packages the detection pipeline into GitHistorian, an operational tool that can help maintainers and auditors spot alterations in repositories of interest (\Cref{sec:hist-githistorian}).
\Cref{sec:hist-related} and \Cref{sec:hist-methodology} first establish prior-work context and the archival comparison pipeline; the chapter closes with discussion (\Cref{sec:hist-discussion}), threats to validity (\Cref{sec:hist-validity}), and a conclusion (\Cref{sec:hist-conclusion}).

\section{Related Work}
\label{sec:hist-related}

This section reviews related work along the two axes this chapter's contribution rests on: detecting and extracting evidence of history alterations (\Cref{subsec:hist-rw-forensics}), and the motivations behind alterations that make them hard to categorise (\Cref{subsec:hist-rw-rewriting}); it closes by positioning this chapter's contribution against that backdrop (\Cref{subsec:hist-rw-positioning}).

\subsection{Detecting and Extracting Evidence of History Alterations}
\label{subsec:hist-rw-forensics}

Without an independent archive to compare against, existing tooling can only \emph{infer} that history was altered, from clues left in the surviving commits, rather than detect it by direct comparison with what was removed.
AST-based approaches attempt to detect file renames with high accuracy by focusing on content structure rather than metadata~\cite{DBLP:conf/apsec/FujimotoHK21}.
Probabilistic methods have been developed to track file movements between versions---addition, removal, and relocation---using nearest-neighbour clone detection, but without capturing metadata alterations~\cite{DBLP:conf/wcre/LavoieKMZ12}.
Neither line of work operates at the scale of an independent archive, and neither can reconstruct commits that have already been removed from the advertised graph.

There are almost no empirical studies of history alterations, in part because of that same difficulty: after a rewrite, the prior history no longer exists and cannot be compared with what remains.
A single \texttt{git commit --amend}, for instance, leaves no trace in the repository of what the commit looked like beforehand; a researcher who only has access to the current state has nothing to compare it against.
The most relevant prior work subjected the Linux kernel and its contributing repositories to continuous snapshots and found that the actual development history of DVCS projects can be irretrievably lost; that study examined alterations occurring across repositories, including rebasing and metadata changes, but did not examine alterations within individual repositories nor identify specific file-level changes at large scale~\cite{DBLP:journals/ese/GermanAH16}.
Empirical mining research has long warned that Git histories are not a neutral ground truth: early MSR work catalogues structural and statistical pitfalls~\cite{bird2009promisesmininggit}, GitHub-scale studies echo that measurement and sampling are fragile under real platform dynamics~\cite{kalliamvakou2014promises}, and a study of timestamp anomalies in a subset of Software Heritage suggests that history alterations affect timestamp reliability and can compromise dataset integrity~\cite{DBLP:journals/ese/FlintCD22}.
Security research complements those caveats with explicit threat models: metadata manipulation can omit commits, misrepresent merge outcomes, and steer developers toward unsafe actions, motivating signed logs of developer-visible repository state~\cite{torresarias2016gitmetadata}.

Large-scale studies of public development increasingly rely on compressed graph representations of version-control data to make global queries feasible~\cite{saner-2020-swh-graph,DBLP:conf/msr/PietriSZ20}.
Software Heritage has been positioned both as a preservation effort and as a dataset for provenance and reproducibility research because it retains objects and snapshots across upstream churn~\cite{swhcacm2018,cise-2020-doi}.
The altered-histories methodology exploits that property directly: consecutive snapshots of the same origin are differenced, root-cause commits are identified, and alterations are classified along metadata-versus-directory axes.
The stance is forensic (compare two observed worlds) rather than inline cryptographic~\cite{torresarias2016gitmetadata}, because Git itself does not provide append-only project history to downstream clients.

\subsection{Motivations for History Rewriting and the Difficulty of Categorising It}
\label{subsec:hist-rw-rewriting}

Prior work and reference material largely treat history editing as a normal, sanctioned maintenance practice rather than a phenomenon requiring measurement.
The Pro Git book presents history rewriting as a ``great thing'' and lists the actions one can perform \emph{before} sharing development history---modifying files, changing commit messages, and reordering, squashing, splitting, or removing commits---while warning that shared branches require coordination, precisely because rewrites change commit identities and disturb collaboration invariants~\cite{ProGit2014}.
Developers use alterations to maintain ``clean'' histories that ease code review and project management~\cite{bird2009promisesmininggit, DBLP:conf/issre/JiCYM20}, and tooling actively encourages squashing to improve clarity~\cite{DBLP:conf/icse/GousiosPD14, DBLP:conf/scam/PaixaoM19, DBLP:journals/tvcg/KimKJKSKS21}.
That legitimacy is precisely what complicates large-scale analysis: the taxonomy developed in this chapter (\Cref{sec:hist-taxonomy}) must separate such routine hygiene from the governance-sensitive cases examined later, a distinction prior workflow-taxonomy research on rebasing does not need to make~\cite{DBLP:conf/scam/PaixaoM19, DBLP:journals/infsof/RiosEE22}.
History rewriting is also the standard remedy for erroneously committed sensitive information, a use case studied directly in Chapter~\ref{chap:secrets}.

\subsection{Research Gap and Positioning of This Chapter}
\label{subsec:hist-rw-positioning}

To our knowledge, no prior work had combined the two ingredients this characterisation requires: a way to \emph{extract} evidence of alterations from public repositories at scale, which Software Heritage's archive makes possible, and a way to \emph{analyse} what kind of alteration occurred once detected, which requires a taxonomy that did not previously exist for this phenomenon.
This chapter fills that gap: it provides a taxonomy for categorising alterations, digs into two recurrent patterns (secret suppression and retroactive license changes), and provides an automated tool to help developers detect alterations in repositories of interest~\cite{DBLP:conf/kbse/RapaportPTZ25,replication-package-altered-histories}.
The study establishes that public version-control state should not be assumed permanently stable once a repository is published.

Relative to the tag study of Chapter~\ref{chap:tags}, this chapter widens the lens from a single reference type to reachable commit histories: rather than asking whether a tag still resolves to the same commit, it asks whether commits once advertised on public branches disappear or are replaced across successive archival observations, and it foregrounds consequences for forensics and repository evolution rather than dependency resolution~\cite{rapaport2026tagalterations}.

\section{Methodology}
\label{sec:hist-methodology}

This study shares the observational stance established in full in Chapter~\ref{chap:tags} but reaches commit histories rather than tag targets.
The four-phase structure is also shared on some level: collect a large corpus of periodically crawled repositories, detect alterations by snapshot comparison, categorise them with a taxonomy, and analyse recurrent file-level patterns through targeted case studies.
To avoid repeating that shared structure verbatim, the rest of this section focuses on what changes when the object of comparison becomes commit histories rather than tag targets: the unit of differencing, the corpus and filters, and a taxonomy and case studies built specifically for commit-level rewrites.
That shift in object is what makes the two studies' findings substantively different rather than variations on the same result: commit-level comparison exposes file- and directory-content forensics, the security-relevant case studies of \Cref{sec:hist-security} (secret suppression, retroactive license changes), and the GitHistorian tool of \Cref{sec:hist-githistorian}, none of which have a tag-level analogue.

\subsection{From Tag Targets to Reachable Commit Histories}
\label{subsec:hist-differences}

The tag study compares the resolved targets of tag references; this study compares the set of commits reachable from public branch heads.
Because Git identifies each commit recursively over its content, metadata, and parent identifiers (\Cref{sec:bg-git-model}), modifying any part of a commit changes its identifier and---through the snowball effect---the identifiers of every commit that transitively references it.
Figure~\ref{fig:hist-amend-rebase} shows the two common rewrite operations.
In panel~(a), the user amends commit \texttt{2}, producing the new commit \texttt{4}; commit \texttt{2} disappears from the graph and, since commit \texttt{3} is based on commit \texttt{2}, commit \texttt{3} is \emph{also} modified and becomes commit \texttt{5}.
If commit \texttt{2} had one million commits transitively referencing it instead of one, they would \emph{all} disappear from the repository and be replaced by new ones.
In panel~(b), a rebase merges commits \texttt{c} and \texttt{b} into a single new commit \texttt{d}; in other use cases, a rebase can also split commits or attach them to different previous commits.

\begin{figure}[h]
\centering
\includegraphics[width=0.92\linewidth]{git_amend.drawio.pdf}
\caption{Common scenarios leading to history alterations in Git: (a) commit amendment, (b) interactive ``rebase''. Grayed out commits are no longer reachable after alteration. Arrows point from newer to older commits.}
\label{fig:hist-amend-rebase}
\end{figure}

In both cases the old commits disappear from the advertised graph and new ones replace them, while Software Heritage still reaches the old commits from the previous snapshot of the same origin, as shown in figure~\ref{fig:swh-without-time}.

\begin{figure}[h]
  \centering
\includegraphics[width=0.92\linewidth]{swh-without-time.drawio.pdf}
  \caption{Repository state before and after \texttt{\small git commit --amend} of \Cref{fig:hist-amend-rebase}(a).
    Both histories are preserved by \SWH and share common commits.
  }
  \label{fig:swh-without-time}
\end{figure} 

This yields the operational definition used throughout: a repository $R$ \emph{underwent a history alteration} if there exist two subsequent snapshots $S_1, S_2$ of $R$ such that at least one commit present in $S_1$ is missing from $S_2$.
An \emph{altered commit} is one present in an earlier snapshot but missing from a later one.
Because of the snowball effect, a single alteration can affect many commits, so for each alteration the analysis identifies one or more \emph{root-cause commits}: the oldest commits, in topological order, among all altered commits---those are unequivocally directly altered rather than changed as a consequence of an earlier modification.

\subsection{Data Collection and Repository Selection}
\label{subsec:hist-data-collection}

The corpus starts from the Software Heritage graph dataset, timestamp \FullDataset, which contains the VCS history of more than \SWHOriginShort repositories and packages (``origins'')~\cite{DBLP:conf/msr/PietriSZ20}.
To avoid under-estimating the ratio of repositories that witness alterations, origins visited fewer than two times are excluded (\OriginsOneVisitOrLessShort), as are non-Git origins (about \num{19000}), leaving \OriginsTwoVisitsOrMoreShort origins where evidence of history alterations can potentially be found.
As in the tag study, popularity-dependent analyses that relate alteration likelihood to repository visibility use GitHub star counts made available by Software Heritage and are restricted to GitHub repositories, while all other analyses use every archived origin regardless of platform; star counts can be inflated but still indicate visibility and scrutiny-worthiness~\cite{DBLP:journals/corr/abs-2412-13459}.

\subsection{Detection Framework and Branch Unification}
\label{subsec:hist-detection}

For each origin with at least two visits, and for each Git branch in them, all pairs of successive snapshots are compared to identify commits from the earlier snapshot missing from the later one (figure~\ref{fig:hist-methodology} step~1); within that set, root-cause commits are extracted per the definition above (figure~\ref{fig:hist-methodology} step~2).
This step is implemented with custom code that mines the compressed graph representation of the Software Heritage archive locally, via its Rust API~\cite{saner-2020-swh-graph}.
The framework counts alterations at repository scale and records branch-level patterns for the distribution analyses in Section~\ref{sec:hist-prevalence}.

\begin{figure}[h]
  \centering
  \includegraphics[width=0.92\linewidth]{flowchart.drawio.pdf}
  \caption{Methodology for detecting and categorizing history alterations.}
\label{fig:hist-methodology}
\end{figure}

Because different branch naming schemes denote the same concept, branches are unified before aggregation (Table~\ref{tab:hist-branch-unification}): \texttt{main} and \texttt{master} become \texttt{main}; \texttt{dev}/\texttt{devel}/\texttt{develop}/\texttt{development} become \texttt{development}; \texttt{pull/<number>/head} becomes \texttt{pull request}; and \texttt{renovate/<anything>} becomes \texttt{renovate}.

\begin{table}[h]
\centering
\caption{Branch-name unification used before aggregating alteration statistics.}
\label{tab:hist-branch-unification}
\begin{tabular}{ll}
\hline
\textbf{Branch names or patterns} & \textbf{Unified name} \\
\hline
\texttt{main}, \texttt{master} & \texttt{main} \\
\texttt{dev}, \texttt{devel}, \texttt{develop}, \texttt{development} & \texttt{development} \\
\texttt{pull/<number>/head} & \texttt{pull request} \\
\texttt{renovate/<anything>} & \texttt{renovate} \\
\hline
\end{tabular}
\end{table}

\subsection{Taxonomy of Altered Commits}
\label{subsec:hist-classification}

Root-cause commits are categorised with a custom taxonomy (Table~\ref{tab:hist-rootcause-categories} and step 3 of figure~\ref{fig:hist-methodology}), building on the closely related taxonomy developed for tag alterations in Chapter~\ref{chap:tags} (\Cref{subsec:hist-rw-positioning}).
It was developed via an iterative mixed method: sampling uncategorised alterations, manually comparing before/after states, capturing the individual data that differed (introducing new categories as needed), applying all categories to the whole dataset, and repeating until everything was categorised.
The categories are:
\begin{itemize}
\item \textbf{Different Branch Name:} content and metadata are unchanged, but the commit is now found only in a different branch than before---the simplest case.
\item \textbf{META:} the commit was replaced by a new commit with strictly identical file and directory content but different metadata; this is refined to indicate which metadata fields changed (author, message, date, committer, committer date, other), and multiple fields can change at once.
\item \textbf{DIR:} the commit's file or directory content was altered or deleted, in part or whole; metadata changes are \emph{not} analysed here because there is no single reference commit to compare against. DIR is refined into \emph{Content Split} (each file that existed in the altered commit is still present but split across several other commits; exclusive of the other DIR subcategories), \emph{File Modified} (at least one file modified), and \emph{File Removed} (at least one file removed and not found elsewhere).
\end{itemize}
To distinguish these, the analysis looks in the newest snapshot at commits starting from the children of the parents of the root-cause commit and proceeds for ten generations (commit-tree depth).
If a file path cannot be found after ten generations of descendants---or if there are no more descendants---the file is considered removed.
The horizon of ten generations is a heuristic trade-off between analysis time and the diminishing likelihood that a file reappearing far from a change is genuinely the same file.

\begin{table}[h]
\centering
\caption{Root-cause commits by category. Metadata fields within META are not mutually exclusive.}
\label{tab:hist-rootcause-categories}
\begin{tabular}{llr}
\hline
\textbf{Category} & \textbf{Field} & \textbf{\#Commits} \\
\hline
\multirow{6}{*}{META (\METARootCauseCommits)}
 & Author & \CategAuthor \\
 & Message & \CategMessage \\
 & Date & \CategDate \\
 & Committer & \CategCommitter \\
 & Committer Date & \CategCommitterDate \\
 & Other & \CategOther \\
\hline
\multirow{3}{*}{DIR (\DIRRootCauseCommits)}
 & File Modified & \CategFileModified \\
 & File Removed & \CategFileRemoved \\
 & Content Split & \CategContentSplit \\
\hline
Different Branch Name & --- & \DBNRootCauseCommits \\
\hline
\end{tabular}
\end{table}

\subsection{Case-Study Design}
\label{subsec:hist-casestudy-design}

Among DIR alterations, two recurring patterns carry the clearest security and legal stakes---as opposed to routine content changes such as dependency version bumps---and are therefore selected for in-depth case studies whose results are reported later, in Section~\ref{sec:hist-security}.
Their detection criteria are fixed here, alongside the taxonomy, because both must be defined before the corpus is processed.
The first, \emph{secret suppression}, focuses on files removed (category: File Removed) that should never have been shared via a public repository, detected by popular filenames used to store private keys or passwords, such as \texttt{id\_rsa} (without the \texttt{.pub} extension) or \texttt{credentials.txt}; the full list ships with the replication package~\cite{replication-package-altered-histories}.
The second, \emph{retroactive license changes}, identifies modified license files (category: File Modified) on the main branch (after unification) whose names contain \texttt{LICENSE} or \texttt{LICENCE} (case-insensitive).
ScanCode, a state-of-the-art license-detection tool~\cite{DBLP:journals/computer/Ombredanne20, scancode-home}, is run on the files before and after alteration, keeping only the highest-confidence result above \SI{90}{\percent}.
The resulting before/after license sets are categorised to gauge the magnitude of change: \emph{License Update} (the set of unversioned licenses is unchanged but a versioned license changes, e.g.\ GPL-2~$\rightarrow$~GPL-3), \emph{Partial Change} (some unversioned licenses changed, e.g.\ GPL added, MIT removed), and \emph{Full Change} (all unversioned licenses changed).

\section{Prevalence and Distribution of History Alterations}
\label{sec:hist-prevalence}

This section presents an empirical analysis of history alterations as they currently stand in the wild: applying the detection framework of \Cref{subsec:hist-detection} to the corpus of \Cref{subsec:hist-data-collection}, it reports how often they occur and how they distribute across branches (\Cref{subsec:hist-dist-branches}) and repositories of different popularity (\Cref{subsec:hist-affected-projects}).

\subsection{Frequency Across Repositories}
\label{subsec:hist-freq}

Across the entire corpus, we observed \AlteredCommitsShort alterations spread across all analysed branches, corresponding to \GitRootCauseCommits (\GitRootCauseCommitsShort) identified root-cause commits.
In terms of repository count, \GitOriginsRootCauseCommits (\GitOriginsRootCauseCommitsShort) repositories contain at least one altered commit, approximately \PercentOriginsAltered of all examined repositories.
The event count is orders of magnitude larger than the repository count because a single maintenance campaign can touch many commits and edges, so readers should not equate ``events'' with ``projects''.
The proportion is modest in relative terms but not negligible in absolute ones: users of more than one million repositories can---and, assuming future stability, will---use public Git repositories that underwent history alterations during their lifetime.
At the same time, the low percentage indicates that history rewriting remains an exception rather than standard practice, suggesting that common Git usage adheres to the expectation of immutable histories.

\subsection{Distribution Across Branches}
\label{subsec:hist-dist-branches}

Alterations are not confined to obscure sandbox branches.
After the unification of Table~\ref{tab:hist-branch-unification}, the top branch categories by share of alteration events are listed in figure~\ref{fig:hist-branch-distribution}: pull-request branches dominate (\HistShareAlterationsPullBranch), reflecting interactive rebasing and commit squashing to clean up history before integration.

However, \texttt{main} (after unifying \texttt{main}/\texttt{master}) alone accounts for \HistShareAlterationsMainBranch of events---ranking second and rejecting the myth that rewriting is only a pull-request hygiene issue.

The third category, \texttt{renovate} (\HistShareAlterationsRenovateBranch), reflects automated dependency-update bots that rewrite commits as proposals evolve before merging, while \texttt{development} branches account for \HistShareAlterationsDevelopmentBranch, ranking fifth overall.
Even though roughly half of all alterations follow a known practice, alterations on main branches compromise fundamental version-control assumptions, potentially breaking downstream dependencies, invalidating signed commits, and disrupting reproducible builds~\cite{10.1145/3641525.3663622}.
The roughly \SI{12}{\percent} alteration rate for main branches suggests that a substantial subset of repositories may pose reliability risks for users who depend on stable commit references.

\begin{figure}[h]
    \includegraphics[width=0.92\columnwidth]{distrib_branch.pdf}
    \centering
    \caption{Most impacted branches by Altered Histories ($> 0.1$\%)}
\label{fig:hist-branch-distribution}
\end{figure}

\subsection{Distribution Across Repositories}
\label{subsec:hist-affected-projects}

Joining the \GitOriginsRootCauseCommitsShort altered repositories with their GitHub star counts (after removing \OriginsUnkownStars origins with unknown star counts) gives the following distribution shown in figure~\ref{fig:hist-popularity-proportion}: 0--1 star, \OriginsZeroOneStar origins (\PercentOriginsZeroOneStar); 2--9 stars, \OriginsTwoNineStars (\PercentOriginsTwoNineStars); 10--99 stars, \OriginsTenNinetyNineStars (\PercentTenNinetyNineStars); 100--999 stars, \OriginsHundredNineHundredNinetyNineStars (\PercentHundredNineHundredNinetyNineStars); and 1000+ stars, \OriginsThousandPlusStars (\PercentThousandPlusStars).

The proportion of origins with a detected alteration in each bucket reveals a pronounced relationship between popularity and alteration likelihood: \HistAltRateZeroOneStar for 0--1 star, rising to \HistAltRateTwoNineStars for 2--9, \HistAltRateTenNinetyNineStars for 10--99, and stabilising around \HistAltRateHundredNineNineNineStars for 100--999 and \HistAltRateThousandPlusStars for 1000+.
It is therefore \emph{not} the case that less popular repositories experience more rewrites; rather, as repositories gain visibility, the likelihood of history modification increases substantially, and even mature, popular repositories (100+ stars) experience a non-negligible amount of alteration.
Within very popular repositories (100+ stars), alterations happen mostly on pull-request branches (\HistPopularPullBranchShare) but remain non-negligible on main branches (\HistPopularMainBranchShare), with implications for dependency management, reproducible builds, and the reliability assumptions developers make when incorporating external repositories.

\begin{figure}[t]
  \includegraphics[width=0.92\columnwidth]{proportion_res_stars.pdf}
  \centering
  \caption{Proportion of Repositories with Altered Histories per Star Range}
\label{fig:hist-popularity-proportion}
\end{figure}

\section{Taxonomy of History Alterations}
\label{sec:hist-taxonomy}

Where \Cref{sec:hist-prevalence} established how often and where alterations occur, this section turns from extraction to analysis: applying the taxonomy of Table~\ref{tab:hist-rootcause-categories} to the \GitRootCauseCommitsShort identified root-cause commits gives the category distribution summarised in figure~\ref{fig:hist-category-distribution}, reported for the whole corpus and for the subset of repositories with 1000+ GitHub stars.
Beyond this distribution, the taxonomy is what makes the rest of the chapter tractable: it is what lets \Cref{sec:hist-security} single out the two DIR patterns---secret suppression and retroactive license changes---that carry the clearest security and legal stakes, rather than treating every alteration as equally concerning.

\begin{figure}[b]
  \includegraphics[width=0.92\columnwidth]{distrib_categ.pdf}
  \centering
  \caption{Distribution of Categories: All vs. Popular Repositories}
\label{fig:hist-category-distribution}
\end{figure}

\subsection{Metadata-Only Alterations (META)}
\label{subsec:hist-tax-metadata}

Among the root-cause commits, \METARootCauseCommitsShort instances (\PercentMETARootCauseCommits) exhibit modifications only to commit metadata, leaving files and directories unaltered.
The committer date is the most frequently modified attribute, across both the complete dataset and the 1000+ star subset, and metadata-modification patterns are uniform regardless of popularity.

A co-occurrence analysis over the metadata fields of Table~\ref{tab:hist-rootcause-categories} shows that when two fields change at once (figure~\ref{heatmap2}), the most frequent pair is the author date together with the committer date---changing one without the other would be surprising.
That same pair reappears at the core of the most frequent triplet (figure~\ref{heatmap3}): occurring \HistTripletMessageDatesShare of the time in the global dataset, it is a commit-message change accompanied by changes to the author and committer dates, the expected behaviour when editing a commit message.
Such events still break naive audit assumptions: any workflow that logs ``who landed commit \texttt{abc1234}'' without anchoring external evidence can be retroactively contradicted.

\begin{figure}[t!]
  \begin{minipage}{3cm}
      \subtop[]{\includegraphics[height=0.3\textheight]{heatmap2.png}\label{heatmap2}}
  \end{minipage}
  \hspace{6cm}
  \begin{minipage}{4cm}
      \subtop[]{\includegraphics[height=0.3\textheight]{heatmap3.png}\label{heatmap3}}
  \end{minipage}
  \caption{Co-occurrence analysis of metadata alteration (a)~Pair (b)~Triplet\TODO{Add mask for pair in triplet and 100\% in pairs}}
  \label{fig:heatmap}
  \end{figure}

\subsection{File-Modifying Alterations (DIR)}
\label{subsec:hist-tax-file}

The taxonomy categorises \DIRRootCauseCommitsShort altered commits by changes to file or directory content; combined with structural changes, file and directory alterations account for \PercentDIRRootCauseCommits of root-cause commits.
The three DIR subcategories carry different stakes rather than being a bookkeeping detail: \emph{File Modified} alterations preserve every path and are the closest analogue to silent drift, \emph{File Removed} alterations erase content or rename it while modifying (\Cref{sec:hist-security}).
We identify only remanes that are keeping the exact same content and rename the filename exclusively, we categorise them as \emph{File Modified}.
Finally \emph{Content Split} alterations scatter a file's history across several descendant commits, complicating even archival recovery.

\emph{File Modified} dominates in both datasets (all vs. popular repositories, figure~\ref{fig:hist-category-distribution}), exceeding \SI{60}{\percent} of occurrences, and is more than four times as frequent as file removal: a developer who alters a file in history is roughly four times more likely to modify it than to delete, rename, or move it.
These events are the closest analogue to ``silent code drift'' for consumers who track branch tips: the branch name is stable while the underlying tree-content lineage changes.

\subsection{Structural and Reference-Level Changes}
\label{subsec:hist-tax-structural}

The \emph{Different Branch Name} subcategory accounts for \DBNRootCauseCommitsShort commits whose content and metadata are unchanged but that are now found only in a different branch.
The most plausible explanation is the merge of so-called feature branches that are later removed; this is still problematic for direct users of those branches, though understandable when the branches are explicitly documented as volatile.
Grouping these phenomena alongside file and metadata signals ensures that prevalence tables do not conflate ``tree stable'' with ``history stable''.

\subsection{Other Observed Alteration Patterns}
\label{subsec:hist-tax-other}

Automation-heavy namespaces---for example, dependency-update bots such as Renovate and Dependabot---contribute recurring rewrite shapes that are operationally benign but still appear as alterations in the archive.
The taxonomy therefore supports separating human-maintained default branches from bot-owned integration lines when interpreting the branch tables in Section~\ref{sec:hist-prevalence}.

\section{Security-Relevant Characteristics of Altered Histories}
\label{sec:hist-security}

This section moves from category counts to concrete security stakes: it examines which files are actually targeted (\Cref{subsec:hist-sec-files}) and then follows the two DIR patterns selected in \Cref{subsec:hist-casestudy-design}---retroactive license changes (\Cref{subsec:hist-sec-license}) and secret removal (\Cref{subsec:hist-sec-secrets})---in depth.

\subsection{Recurrent File Patterns}
\label{subsec:hist-sec-files}

Because file alterations are the most common form of rewriting, an obvious question is which files are modified or removed.
Extracting the filename part of all paths involved in DIR alterations yields the top-20 most commonly altered filenames in Table~\ref{tab:hist-top-filenames}.
The most frequently altered files are predominantly configuration and build-related artifacts: \texttt{default.nix}, \texttt{package.json}, \texttt{Makefile}, \texttt{pom.xml}, \texttt{CMakeLists.txt}, and \texttt{APKBUILD} represent different build systems and package managers, suggesting that dependency updates, version bumps, and build-configuration changes are primary motivations for rewriting.
The high frequency of \texttt{default.nix}, \texttt{package.json}, and \texttt{meta.yaml} matches the workflows of automated dependency tools (Renovate, Dependabot, Nix update bots) that modify these files and then rewrite commit histories.
Documentation files (\texttt{README.md}, \texttt{CHANGELOG.md}, \texttt{index.md}) and ecosystem entry points (\texttt{index.js}, \texttt{index.html}, \texttt{index.ts}, \texttt{\_\_init\_\_.py}) also appear frequently, while \texttt{LICENSE} (\#20) and \texttt{metadata.xml} (\#13) indicate that legal-compliance and metadata corrections constitute a notable category---motivating the license case study below.

\begin{table}[h]
\centering
\caption{Top-20 most altered file names.}
\label{tab:hist-top-filenames}
\small
\begin{tabular}{rlr@{\hskip 2em}rlr}
\hline
\textbf{\#} & \textbf{File name} & \textbf{Count} & \textbf{\#} & \textbf{File name} & \textbf{Count} \\
\hline
1 & \texttt{default.nix} & \num{13401935} & 11 & \texttt{CHANGELOG.md} & \num{5001971} \\
2 & \texttt{package.json} & \num{12594717} & 12 & \texttt{\_\_init\_\_.py} & \num{4859517} \\
3 & \texttt{Makefile} & \num{10615195} & 13 & \texttt{metadata.xml} & \num{4477947} \\
4 & \texttt{README.md} & \num{9969599} & 14 & \texttt{CMakeLists.txt} & \num{4380963} \\
5 & \texttt{index.js} & \num{9129894} & 15 & \texttt{package.py} & \num{4336900} \\
6 & \texttt{Manifest} & \num{8891876} & 16 & \texttt{APKBUILD} & \num{3375064} \\
7 & \texttt{pom.xml} & \num{7080355} & 17 & \texttt{template} & \num{3122049} \\
8 & \texttt{index.html} & \num{6297643} & 18 & \texttt{BUILD} & \num{3018131} \\
9 & \texttt{meta.yaml} & \num{5425074} & 19 & \texttt{index.ts} & \num{2764520} \\
10 & \texttt{index.md} & \num{5372929} & 20 & \texttt{LICENSE} & \num{2762147} \\
\hline
\end{tabular}
\end{table}

History alterations predominantly target infrastructure, configuration, and metadata files rather than core application logic: maintenance activities and automated tooling, not attempts to hide application behaviour, are the primary drivers of the observed alterations.
\subsection{Case Study: Retroactive License Changes}
\label{subsec:hist-sec-license}

The first case study focuses on retroactive license modifications, where developers alter history to change the content of license-denoting files.
We identified \LicenseAll altered license files across the main branches of the dataset, spanning \LicenseOrigins repositories, \LicenseOriginsThousandPlusStars of which have 1000+ GitHub stars.
Of these, \LicenseRemoved instances involved renaming, relocation, or deletion, while \LicenseModified represented content modifications.

Using ScanCode, we identified the before/after license sets for \LicenseModifiedBothScanned alterations: \LicenseUpdate are relatively minor license updates (at most a change in license version, with copyright and author edits a common occurrence in this category), \LicenseFullChange represent full license changes (including transitions from permissive licenses such as MIT to restrictive ones such as GPL and vice versa), and \LicensePartialChange are partial changes, the last attributable to the limited number of repositories (\LicenseSeveralFamilies alterations) with multiple license families.

These findings are concerning for downstream users: while most open-source licenses are legally irrevocable, benefiting from the granted rights requires practical access to a version of the software released under that license.
Retroactively changing version history can make old versions, under an old license, disappear from public circulation---inhibiting previously available rights, an especially acute concern for industrial users who expect legal stability.

\subsection{Case Study: Removing Secrets from History}
\label{subsec:hist-sec-secrets}

Across all repositories, we identified \SecretsRemovedAllShort history alterations categorised as File Removed that involve files whose names commonly denote private information (private keys, certificates, passwords), spanning \SecretsOriginsAllShort distinct repositories.
Generic naming patterns are the most frequent: files containing ``key'' and ``secret'' account for \SecretsRemovedAllKeyShort and \SecretsRemovedAllSecretShort altered files respectively.
Such removals concentrate where rewriting is routine---\PercentSecretsPullBranch of them occur on pull-request branches and \PercentSecretsMainBranch on main branches overall, rising to \PercentSecretsThousandStarsPullBranch on pull-request branches within 1000+ star repositories.
Manual examination of a representative random sample confirmed that removed files contained sensitive information, in files such as \texttt{secret.yaml} (authentication credentials) and \texttt{samlKey.jks} (keys for deploying production environments); filename heuristics alone, however, cannot exclude removals driven by refactoring, false-positive path matches, or rewrites whose primary purpose was unrelated to credential exposure.
SSH private keys are less common---probably because platforms such as GitHub can block pushes containing them---but are not absent: \SecretsRemovedIdRsa removed files contained RSA private keys, spanning \SecretsRemovedIdRsaOrigins origins.

These findings show that despite attempts to remediate exposure through history alteration, sensitive information remains recoverable through archival history digging: removing a secret from history is only partial remediation, and complete restoration requires regenerating or rotating all exposed credentials---at which point there is arguably little point in rewriting history afterwards.
It is important to note that this methodology only detects alterations occurring between captured repository snapshots: the results correspond either to cases where a snapshot happened to be captured immediately after the secret was committed, or to instances where significant time elapsed between exposure and the remediation attempt.

\section{GitHistorian: An Operationalization of History Alteration Detection}
\label{sec:hist-githistorian}

Archival counts establish prevalence; practitioners still need \emph{local} explanations---what disappeared, what replaced it, and which branch heads witnessed the change.
History alterations otherwise go unnoticed except by downstream users who happen to have retrieved a repository version before the alteration and then pulled after it with branch tracking in place---conditions few users satisfy, since CI/CD pipelines often perform a fresh clone, making detection impossible.
Getting that local explanation directly from Software Heritage is not a realistic option for most practitioners: the archive is built for large-scale preservation and graph queries, not for on-demand, repository-by-repository auditing, and its Rust API and compressed graph format carry a steep learning curve that maintainers auditing a single dependency should not have to climb.
GitHistorian closes that gap: it is a tool that detects, describes, and (if desired) helps avoid repositories that witnessed history alteration, packaging the detection logic of \Cref{sec:hist-methodology} behind an interface that does not require Software Heritage expertise.

\subsection{Design}
\label{subsec:hist-gh-goals}

The design satisfies three core requirements: scalability for analysing thousands of repositories, accuracy in detecting various alteration types, and usability for both interactive analysis and automated monitoring.
The architecture has four components: (1)~a PostgreSQL database storing preprocessed history-alteration data derived from Software Heritage snapshots; (2)~a command-line interface for interactive analysis; (3)~a caching mechanism to optimise repeated queries; and (4)~Git-hook integration for continuous monitoring.
The local database is required because Software Heritage does not currently expose a remote API answering the queries needed to detect history alterations; this requirement could be lifted in the future if the archive operators or a third party provided such a service over the periodically published datasets.
The detection mechanism queries the database for known alterations and supports branch-specific analysis (main branches only, development branches only, or all branches), and for automated monitoring it installs Git hooks that trigger checks after repository updates (\texttt{post-merge}) and branch switches (\texttt{post-checkout}).

\subsection{Implementation and CLI}
\label{subsec:hist-gh-workflow}

GitHistorian is implemented in Rust, using the \texttt{sqlx} crate for type-safe database interactions and \texttt{tokio} for asynchronous operations, enabling efficient handling of datasets containing billions of alteration records.
The main functionality is exposed through CLI commands:
the \texttt{load} command ingests the dataset into the local database with parallel processing;
the \texttt{check} command audits a specific repository, querying the local database and formatting results for human consumption (with optional restriction to main or development branches);
the \texttt{attach} command installs Git hooks that call \texttt{check-cached}, a cached variant that keys results by repository URL, target branch, and dataset version to avoid redundant queries.
Cache validity is determined by dataset-version comparison, so results remain accurate as new alteration data becomes available while repeated queries return immediately.
Hook scripts include contextual information about the triggering event (pull, merge, or checkout) and provide clear user feedback about the analysis results; the implementation is available as part of the replication package~\cite{replication-package-altered-histories}.

\begin{figure}[h]
  \centering
  \begin{verbatim}
  $ git-historian check https://github.com/example/project \
      --branch main --verbose
  Connected to the database!
  Found 2 altered history records for
    'https://github.com/example/project'
  Record #1:
    Branch Name: refs/heads/master
    Altered Commit: swh:1:rev:a1b2c3d4e5f6789...
    Snapshot Destination: swh:1:snp:abcd1234...
    Sub Category: FileModified
  Record #2:
    Branch Name: refs/heads/dev
    Altered Commit: swh:1:rev:f6e5d4c3b2a1098...
    Sub Category: CommitterDate

  $ git-historian attach . --branch main --verbose
  git-historian successfully attached to repository
    Hooks installed:
    - post-merge (triggered after git pull)
    - post-checkout (triggered after git checkout)
  \end{verbatim}
  \mycaption[GitHistorian CLI session]{%
  Representative GitHistorian session: auditing a repository for alterations and attaching Git hooks for continuous monitoring. SWHIDs are illustrative.%
  }
  \label{fig:hist-gh-transcript}
  \end{figure}

\subsection{Example Outputs and Usage Scenarios}
\label{subsec:hist-gh-examples}

A maintainer can check whether a repository's history was altered, then attach the tool for ongoing monitoring, as illustrated in Figure~\ref{fig:hist-gh-transcript}.
Documented scenarios include pre-merge auditing of third-party forks, periodic organisation-wide scans, and CI gates that alert when unexpected rewriting appears.
The verbose output reports each altered commit's branch, SWHID, snapshot destination, and subcategory, and surfaces modified file paths---information that can be crucial for security auditing (for example, when a security-related file such as \texttt{src/security/auth.py} was modified in an altered commit).

\subsection{Practical Usefulness and Limitations}
\label{subsec:hist-gh-usefulness}

GitHistorian shifts the cost of \emph{spot checks} from forensic expertise to tooling, but it inherits archival limits: if Software Heritage never captured the pre-rewrite state, there is nothing to diff against.
It also does not, by itself, classify intent; it surfaces \emph{that} rewriting occurred and \emph{what} classes of tree or metadata edits accompany it.

\section{Discussion}
\label{sec:hist-discussion}

The measurements of \Cref{sec:hist-prevalence,sec:hist-taxonomy,sec:hist-security} establish \emph{that} history alterations are widespread and \emph{what} they touch; this section's objective is to argue \emph{why that should change how practitioners, tool builders, and researchers treat commit history}---as a record that can be tampered with, not an immutable ground truth---and what each of those audiences should do differently as a result.

\subsection{Implications for Software Development Practices}
\label{subsec:hist-disc-audit-trail}

The findings challenge conventional wisdom about version-control best practices.
The paradoxical correlation between popularity and alteration frequency suggests that current workflows may be fundamentally at odds with the immutability principle underlying version control.
While alterations on pull-request and dependency-management branches follow accepted practice, their normalisation may create cultural acceptance that inadvertently extends to main branches.
The substantial presence of main-branch alterations is a systemic challenge: beyond immediate impacts on CI/CD systems and semantic versioning, these practices undermine the cryptographic integrity guarantees that many security and compliance frameworks depend upon, so organisations relying on commit signatures or audit trails may unknowingly operate with compromised assurance levels.

\subsection{Security and Compliance Implications}
\label{subsec:hist-disc-benign-concealment}

The secret-removal analysis reveals organisational security-governance issues beyond immediate credential exposure: the prevalence of post-hoc secret removal suggests systematic failures in preventive measures, pre-commit hooks, developer training, and secure workflows.
Organisations discovering secrets in their histories face a false choice between public exposure and historical integrity.
More concerning is the implicit assumption that history cleaning provides adequate remediation; the detection capabilities here demonstrate that ``deleted'' secrets remain discoverable through historical analysis, yet many organisations may believe their incident response is complete after a rewrite.
The licensing case study exposes novel legal risks: retroactive license changes create ambiguous precedents where different project versions operate under potentially conflicting terms, introducing intellectual-property uncertainty that traditional license scanning cannot detect.

\subsection{Tool, Ecosystem, and Research Implications}
\label{subsec:hist-disc-auditing}

This subsection's angle is downstream of the previous one: rather than the harms of individual alterations, it asks what the \emph{normalisation} of rewriting does to the tools that build on commit history and to the research that mines it.
The automation-driven nature of many alterations reveals a disconnect between tool capabilities and user understanding: modern dependency-management systems normalise history rewriting as an implementation detail, and this normalisation effect may contribute to the casual application of history alteration techniques in contexts where they are inappropriate.
Security tooling faces the challenge of distinguishing legitimate maintenance from suspicious activity, requiring context analysis beyond simple pattern matching.
The research ecosystem faces a methodological challenge as well: significant portions of repository-based studies may analyse sanitised rather than authentic development data, with implications for the reproducibility of software-engineering research that relies on commit-based metrics and temporal analysis.
The concrete proposal that follows from this is that security tooling and mining pipelines alike should not treat a repository's current branch tip as ground truth by default: both should budget for an explicit archival-comparison step---using an independent archive such as Software Heritage, or GitHistorian's packaging of it (\Cref{sec:hist-githistorian})---before trusting commit-based metrics or auditing conclusions.

\subsection{Broader Research and Community Implications}
\label{subsec:hist-disc-research}

The prevalence of history modifications raises a methodological concern for the software-engineering research community that extends beyond the security applications discussed above.
Repository-based studies---mining commit timelines, measuring development velocity, training defect prediction models, or constructing contributor networks---implicitly treat the commit graph as a faithful record of development history.
The measurements reported here suggest that a non-trivial portion of that data may reflect sanitised or restructured versions of the original development sequence rather than authentic event logs; researchers who mine repositories without anchoring to independent archival snapshots inherit this uncertainty.
This blind spot is especially consequential for longitudinal studies and for research that relies on commit-based metrics and temporal analysis, where the order, authorship, and content of commits are all treated as ground truth.
For the broader open-source ecosystem, the findings raise questions about trust, reproducibility, and the long-term integrity of collaborative development platforms.
Developing community standards, audit practices, and tooling to surface and communicate the presence of history alterations represents an important area for future investment.
Platform operators could support this audit posture with append-only logs of ref updates, including force-pushes that overwrite branch tips: each entry would record the branch name, the author of the rewrite, the previous commit hash, and the new tip hash.
Such logs would make post-hoc remediation attempts visible to auditors even when file content was deliberately removed from history, because the evidence does not need to blobs---only the commit graph movement.

\subsection{Supply-Chain Attack Obfuscation}
\label{subsec:hist-disc-supplychain}

This closes the loop with the supply-chain framing set out in \Cref{sec:hist-intro}: a moved tag misdirects \emph{which} commit is fetched, but a history alteration changes \emph{what} that commit contains, and it is this second, content-level capability that an adversary can weaponise.
History alterations present a concerning attack vector: an adversary could inject malicious code while using history rewriting to erase evidence, making detection through conventional auditing extremely difficult.
The documented normalisation of rewriting creates cover for malicious actors to disguise backdoors as legitimate maintenance, and the scale observed---affecting millions of repositories---suggests a significant blind spot, since traditional security measures focus on current repository states and may miss attacks that exploit the temporal dimension of version-control history.
Supply-chain guidance increasingly stresses pinning and provenance; the measurements here show \emph{why} pinning to a branch tip without a cryptographic object id is fragile, and why independent archives matter when disputes arise about what was public when.

\section{Threats to Validity}
\label{sec:hist-validity}

\subsection{Construct Validity}
\label{subsec:hist-val-bias}

History alteration is, by design, a destructive operation; observing it faithfully without information loss would require a user study on developer machines, dramatically reducing scale.
This work takes the opposite approach---a very large-scale experiment relying on Software Heritage to collect and preserve snapshots that are then compared pairwise---and therefore depends on archival frequency: only alterations visible between successive visits are observable.
The prevalence figures reported in this chapter should be read as a conservative lower bound on the phenomenon: missed or merged events imply undercounting, not inflation, so the true scale of history rewriting and its downstream impact on reproducibility, auditing, and supply-chain assurance may be significantly larger than the archival evidence alone establishes.
Short-lived malicious or accidental alterations reverted between visits are missed, and multiple alterations between two visits of the same repository cannot be distinguished and are merged into one.
Metadata changes are not detectable when file changes co-occur, due to the lack of a single reference commit to compare against; metadata could nonetheless be characterised for more than one million commits, with no claim made about the rest.
The ten-commit search horizon for locating moved content is a deliberate trade-off: content might be found with a larger horizon, but a file reappearing many commits away is debatably the same file.

\subsection{External Validity}
\label{subsec:hist-val-external}

Only Git repositories are considered, and no further generality to other DVCS is claimed.
Software Heritage archives other DVCS as well, so the same dataset could in principle support comparative analysis, but disentangling the specificities of each system (for example, branch naming) would be difficult, and some DVCS support non-destructive history updates that retain previous states.
Finally, Git is also used as a storage backend for independent DVCS such as Jujutsu (\texttt{jj}), where history alterations are used more liberally; these cases were not separated from ``regular'' Git repositories and may affect the analysis, though they are expected to remain quantitatively marginal to date.

\section{Conclusion}
\label{sec:hist-conclusion}

This chapter showed that history rewriting on public Git repositories is a large-scale, archivally observable phenomenon: across \OriginsTwoVisitsOrMoreShort repositories, we found \AlteredCommitsShort alteration events and \GitOriginsRootCauseCommitsShort affected repositories (\PercentOriginsAltered), with branch distributions that implicate default lines rather than only pull-request sandboxes.
The phenomenon is not negligible: downstream users of those repositories might have experienced disruption of their push/pull workflow, or might be currently using repositories that have been tampered with without their knowledge.
Applying a novel taxonomy, most rewriting touches files or directories (\PercentDIRRootCauseCommits of root-cause commits), while a smaller but material share is metadata-only (\PercentMETARootCauseCommits), which matters for any audit process that trusts commit hashes without examining trees.
Two targeted case studies showed that alterations are used to remove private information committed by mistake (for example, private keys) and to retroactively change open-source licenses; GitHistorian demonstrates that the same evidence base can feed operational tooling, not only global statistics---a necessity given that history alterations are, by design, meant to leave no traces.

Future work includes discerning legitimate from malicious alterations and designing automated detection of suspicious patterns, and studying the \emph{timing} of rewrites---distinguishing immediate corrections from long-term exposures, and, in the secret-suppression scenario, early-stage removals (when a project was young and security practices immature) from more worrisome late-stage removals.

The most urgent open question left by the secret-suppression case study is one that this chapter cannot answer from archival evidence alone.
The \SecretsRemovedAllShort removals of files whose names suggested secrets detected across \SecretsOriginsAllShort repositories show that developers \emph{did} push private keys, credential files, and authentication tokens to public repositories, and then rewrote history in an attempt to erase them.
But erasing a secret from the repository does not retract it from independent archives, forks, caches, or mirrors---and, more critically, it does nothing to the credential itself unless the credential is separately rotated.
The secrets identified in this chapter remain recoverable from Software Heritage regardless of the rewrite; the question is whether they also remain \emph{usable}.
Chapter~\ref{chap:secrets} takes up that question directly: it recovers secrets embedded in alteration-affected commits, characterises their types and exposure windows, and establishes a responsible-disclosure protocol to test whether the underlying credentials are still active after the remediation attempt.
